\begin{frame}{Introduction}

    \begin{figure}
        \centering

        \begin{subfigure}{0.48\textwidth}
            \centering
            \includegraphics[width=\linewidth,height=0.55\textheight,keepaspectratio]{ising_example_clustersize.png}
            \caption{Log-log plot of the cluster-size distribution $n_I(s,V)$ versus $s$ for an Ising model in the Ising sector \cite{fang2021percolation}.}
        \end{subfigure}
        \hfill
        \begin{subfigure}{0.48\textwidth}
            \centering
            \includegraphics[width=\linewidth,height=0.55\textheight,keepaspectratio]{earthquake_power_law}
            \caption{Earthquake magnitude distribution in Southern California showing a power-law behavior over 6 decades \cite{bak2002unified}.}
        \end{subfigure}

        \caption{Two different systems in a critical state showing behavior governed by power laws.}
    \end{figure}

    Question: What is the difference between these two?

\end{frame}



\begin{frame}{Self-organized criticality}
    Answer: The Ising model needs a fine-tuned parameter (e.g. temperature) to reach criticality, while earthquakes exhibit this behavior in their own. \\
        $\implies$ Self-Organized Criticality (SOC) \pause

    \begin{itemize}
        \item SOC naturally explains the widespread appearance of fractal structures and inverse-frequency noise in nature \pause
        \item Scale invariance: events of all sizes happen, no single "characteristic" scale
            $\implies$ Probability distributions of events follow power laws (Pareto distribution)
    \end{itemize}

\end{frame}



%\begin{frame}{Power laws}

    % TODO: Image of one or two power laws, e.g. solar flares (see https://arxiv.org/pdf/2212.02678)

%\end{frame}



\begin{frame}{Why sandpiles?}

    \begin{itemize}
        \item The Bak-Tang-Wiesenfeld (BTW) sandpile model is a very simple model to study SOC (\cite{bak1987self})
        \item Analogy: Slowly dropping sand onto a table causes regular avalanches of various sizes
    \end{itemize}

    \begin{figure}
        \centering
        \includegraphics[width=\linewidth,height=0.55\textheight,keepaspectratio]{sandpile_example}
        \caption{Example state of a two dimensional BTW sandpile model.}
    \end{figure}

\end{frame}



\begin{frame}{Goals of our project}

    \begin{itemize}
        \item Simulate a BTW sandpile cellular automaton in $d=2,3,4$ dimensions
        \item Show that the system exhibits self-organized criticality
        \item Confirm the emergence of power laws and compute the scaling exponents
        \item Compare the results for different hyperparameters and dimensions
    \end{itemize}

\end{frame}
